\documentclass[../main.tex]{subfiles}
\begin{document}
%==========================================TIÊU ĐỀ TẬP========================================
\begin{table}[h]
  \begin{adjustwidth}{-1cm}{}
    \begin{tabular}{>{\centering\arraybackslash}p{6cm}|>{\raggedright\arraybackslash}p{12.5cm}}
      \multirow{5}{6cm}
      % Cột bên trái: ÉPISODE và số tập
      {\\[-40pt]
        \flaregothic\fontsize{20pt}{20pt}\selectfont \centering
        \color{eptype}ÉPISODE\color{black}\\
        \burbank\fontsize{72pt}{72pt}\selectfont
        \color{epnum}126\color{black}\\[6pt]
        \cafeteria\fontsize{20pt}{20pt}\selectfont\color{epnum}(10-6)\color{black}
        \\[18pt]
        % \flaregothic\fontsize{14pt}{14pt}\selectfont
        % \color{epattr}BIENVENUE, \uline{\color{Banpire} Mel} !
        % \flaregothic\fontsize{18pt}{18pt}\selectfont
        % \color{epattr}ÉPISODE PERDU
        \color{black}
      }
       &
      % Dòng thứ nhất: tên tập tiếng Nhật
      {\fotskip\fontsize{18pt}{18pt}\selectfont
        \ruby{起}{お}こすな\ruby{危}{き}\ruby{険}{けん}！
      }                                              \\[6pt]
      % Dòng thứ hai: tên tập tiếng Anh
       & {\cafeteria\fontsize{18pt}{18pt}\selectfont Do Not Wake}                           \\[4pt]
      % Dòng thứ ba: tên tập tiếng Việt
       & {\vnmsans\fontsize{18pt}{18pt}\selectfont Bài thơ lục bát về một con chó đang ngủ} \\[8pt]
      % Dòng thứ tư: Nhân vật xuất hiện
       & {\sen{\fontsize{14pt}{14pt}\selectfont Track used: Sunny (\ruby{晴}{は}れ)}}
      \\[6pt]
      % Dòng cuối cùng: Thời gian diễn ra của tập (thường sẽ là ngày phát hành)
       & {\continuum\fontsize{16pt}{16pt}\selectfont Ngày phát hành: 10/10/2021}            \\[18pt]
    \end{tabular}
  \end{adjustwidth}
\end{table}
%=========================================TRUYỆN CHÍNH========================================
\color{black}

\begin{center}
  \uline{(Như tiêu đề của tập, câu chuyện này sẽ được viết dưới dạng thể thơ lúc bát, vì tập này không có thoại, ngoại trừ của A-chan ở cuối tập và ở {\abv Truyện phụ}.)}
\end{center}

\il

\begin{center}
  Tháng Mười trời thắm màu thu,\\
  Một nàng Cún nâu nằm ngủ say lành.\\
  Sấm Âm\footnote{Tức Korone, lấy từ tên Hán-Việt 沁音.} tên gọi nhẹ tanh,\\
  Nằm nơi sàn gạch tưởng như thảm mềm.\\
  Tay chân ưỡn thẳng êm đềm,\\
  Ngáy khò khò nhẹ như thêm khúc đàn.\\
  Giấc mơ lấp lánh mơ màng,\\
  Chiều thu ấm áp ru nàng ngủ ngon.\\

  \img{10-6a}

  Bầu Trời\footnote{Tức Sora, lấy từ cụm Hán-Nôm.} cùng Cửu Viễn\footnote{Tức Kuon, lấy từ tên Hán-Việt 久遠.} tròn,\\
  Vừa qua cửa ngõ đã còn ngỡ ngàng.\\
  Khuyển kia nằm ngủ mơ màng,\\
  Trời cười, Viễn chỉ thở than nhẹ lời.\\
  ``Em nằm đây giữa đất trời\dots''\\
  Chống tay Viễn ngước, miệng cười mà than.\\
  Trời ôm Sấm Âm dịu dàng,\\
  Đặt lên ghế bành êm ái như mơn.\\
  Trời dặn khẽ: ``Hãy nhớ luôn,\\
  Đừng làm em giật, tội con chó này.''

  \begin{figure}[H]
    \centering
    \begin{subfigure}[t]{0.45\textwidth}
      \centering
      \includegraphics[width=8cm]{../../../../Image/Book?/10-6b1.png}
    \end{subfigure}
    \quad
    \begin{subfigure}[t]{0.45\textwidth}
      \centering
      \includegraphics[width=8cm]{../../../../Image/Book?/10-6b2.png}
    \end{subfigure}
  \end{figure}

  Có hai chiếc bóng nấp ngay,\\
  Tiểu Chúc\footnote{Tức Okayu, lấy từ Hán-Việt 小粥.}, nàng Tế\footnote{Tức Matsuri, lấy từ Hán-Việt 祭.} chốn này lặng im.\\
  Thấy nàng Khuyển vẫn im lìm,\\
  Chúc đưa tay vuốt dịu êm bạn mình.\\
  Cờ đuôi vẫy nhẹ rung rinh,\\
  Dường như nàng biết ân tình bạn xưa.\\
  Tế nhìn, mắt sáng như lửa,\\
  Cũng đưa tay vuốt chẳng thua ai nào.\\
  Ấy vậy Khuyển chẳng động vào,\\
  Nằm như thể ngủ vùi vào mộng bay.\\
  Tế thôi, mặt cúi lặng thay,\\
  Viễn nhìn bèn bảo: ``Thấm này lắm chiêu!''\\
  Anh chàng cũng thử làm liều,\\
  Vừa xoa, đuôi Khuyển khua nhiều dữ ghê.\\
  Tế nhìn, lòng chỉ ê chề,\\
  Chúc đành cười nhẹ, xoa khe khẽ đầu.

  \begin{figure}[H]
    \centering
    \begin{subfigure}[t]{0.45\textwidth}
      \centering
      \includegraphics[width=8cm]{../../../../Image/Book?/10-6c1.png}
    \end{subfigure}
    \quad
    \begin{subfigure}[t]{0.45\textwidth}
      \centering
      \includegraphics[width=8cm]{../../../../Image/Book?/10-6c2.png}
    \end{subfigure}
  \end{figure}

  Bỗng đâu lại có thêm người,\\
  Vĩnh Viễn\footnote{Tức Towa, lấy từ Hán-Việt 永遠.} xuất hiện giữa thời khắc vui.\\
  Giơ máy chụp ảnh mấy mươi,\\
  Thêm vài bộ lọc, thú cười lố lăng.\\
  Cửu Viễn dẫu mặt nghiêm trang,\\
  Cũng hùa vài kiểu chẳng màng chính chuyên.\\
  Mai Lộ\footnote{Tức Mel, lấy từ Hán-Việt 梅露.} cũng chẳng đứng yên,\\
  Muốn chung bức ảnh giữ riêng kỷ niệm.\\
  ``Cắn yêu nhé\footnote{Ý nói đến câu ``Kappu\~{}'' của Mel.}'', ``Chúc buổi đêm\footnote{Tương tự như chú thích trước với ``Otsuyappi\~{}'' của Towa.}'',\\
  ``Đưa tay cho tôi\footnote{Tương tự như chú thích trước với ``Yubi Yubi!'' của Korone.}, ``Mặt xem ngu ghê!''\\
  Rồi thì lời chúc Tết về,\\
  Chẳng liên quan cũng hả hê miệng cười.\\
  Khuyển kia vẫn ngủ thảnh thơi,\\
  Chẳng hay mình đã\dots\ lên ngôi trò đùa.\\
  Mấy người cơ hội đủ vừa,\\
  Cười rúc rích đến tận trưa chưa ngừng.\\

  \begin{figure}[H]
    \centering
    \begin{subfigure}[t]{0.45\textwidth}
      \centering
      \includegraphics[width=8cm]{../../../../Image/Book?/10-6d1.png}
    \end{subfigure}
    \quad
    \begin{subfigure}[t]{0.45\textwidth}
      \centering
      \includegraphics[width=8cm]{../../../../Image/Book?/10-6d2.png}
    \end{subfigure}
    \quad
    \begin{subfigure}[t]{0.45\textwidth}
      \centering
      \includegraphics[width=8cm]{../../../../Image/Book?/10-6d3.png}
    \end{subfigure}
    \quad
    \begin{subfigure}[t]{0.45\textwidth}
      \centering
      \includegraphics[width=8cm]{../../../../Image/Book?/10-6d4.png}
    \end{subfigure}
  \end{figure}

  Mã Lâm\footnote{Tức Marine, lấy từ Hán-Việt 瑪琳.}, Tuệ Tinh\footnote{Tức Suisei, lấy từ Hán-Việt 彗星.} bước vào,\\
  Cả hai nhìn Khuyển nằm sao ngủ lì.\\
  Chàng giơ ngón trỏ, khẽ thì,\\
  Miệng cười chúm chím, mắt thì long lanh.\\
  Thêm trò tiêu khiển cho nhanh,\\
  Thuyền Trưởng rón rén đưa bánh tới gần.\\
  Bánh thơm chạm mũi ngập ngừng,\\
  Khuyển ngoạm một miếng, chẳng cần suy tư.\\
  Cả hai ôm miệng cười khù,\\
  Dẫu không một tiếng, vẫn như phá làng.\\
  Mã Lâm đắc ý lấn sang,\\
  Chẳng ngờ bị Khuyển ngoạm ngang vào tay.\\
  Im lìm chẳng dám kêu ai,\\
  Tuệ Tinh đổ mồ hôi dài ướt lưng.\\
  Viễn nhìn, miệng mỉm ngậm ngừng:\\
  ``Chọc chó mà thiếu trình từng - thế thôi.''

  \begin{figure}[H]
    \centering
    \begin{subfigure}[t]{0.45\textwidth}
      \centering
      \includegraphics[width=8cm]{../../../../Image/Book?/10-6e1.png}
    \end{subfigure}
    \quad
    \begin{subfigure}[t]{0.45\textwidth}
      \centering
      \includegraphics[width=8cm]{../../../../Image/Book?/10-6e2.png}
    \end{subfigure}
    \quad
    \begin{subfigure}[t]{0.45\textwidth}
      \centering
      \includegraphics[width=8cm]{../../../../Image/Book?/10-6e3.png}
    \end{subfigure}
  \end{figure}

  Lại thêm lũ mới lẻ loi,\\
  Là Thi Âm\footnote{Tức Shion, lấy từ Hán-Việt 詩音.}, Thỏ, với người Thiên phương\footnote{Tức Kanata, lấy từ chữ Hán 方 của 彼方 (Bỉ Phương).}.\\
  Chàng Viễn nhíu mày một đường,\\
  ``Lũ này lại tính tìm đường nghịch chi?''\\
  Phù thủy chạm khẽ thầm thì,\\
  Khuyển không động đậy, cả khi cựa mình.\\
  Họ cười, Viễn lắc đầu khinh,\\
  ``Khôn ngoan gì kiểu dọa thình lình ai?''\\
  Thi Âm liều lĩnh ra tay,\\
  Vừa cù đã bị đá bay ngang tường.\\
  Viễn cười: ``Anh đã lên giường-\\
  (À không), đã bảo đừng vương tay mà\dots''\\
  Hai người kia mặt nhạt nhòa,\\
  Thiên Sứ cũng thử, bị đà tung lên!\\
  Đập trần treo lửng hờ hênh,\\
  Viễn nhìn hãi quá, lắc nhanh đầu buồn.\\
  Thỏ kia tái mặt như suông,\\
  Thở ra chẳng dám động luôn một lần.\\
  Chàng Viễn lặng lẽ khẽ răn:\\
  ``Thích làm trò ngốc, thì ăn đá liền.''

  \begin{figure}[H]
    \centering
    \begin{subfigure}[t]{0.45\textwidth}
      \centering
      \includegraphics[width=8cm]{../../../../Image/Book?/10-6f1.png}
    \end{subfigure}
    \quad
    \begin{subfigure}[t]{0.45\textwidth}
      \centering
      \includegraphics[width=8cm]{../../../../Image/Book?/10-6f2.png}
    \end{subfigure}
    \quad
    \begin{subfigure}[t]{0.45\textwidth}
      \centering
      \includegraphics[width=8cm]{../../../../Image/Book?/10-6f3.png}
    \end{subfigure}
    \quad
    \begin{subfigure}[t]{0.45\textwidth}
      \centering
      \includegraphics[width=8cm]{../../../../Image/Book?/10-6f4.png}
    \end{subfigure}
  \end{figure}

  Bạn A\footnote{Đọc tắt của A-chan.} xuất hiện, hồn nhiên,\\
  Hét vang: ``\textbf{Chào nhé mọi miền gần xa!}''\\
  Khắc Lạp\footnote{Tức Pekora, lấy từ Hán-Việt (佩)克拉.} lập tức xanh da,\\
  Viễn giật bắn cả người ra phía sau.\\
  Nghe vang động, Khuyển nhíu đầu,\\
  Dần dần tỉnh giấc, ngó mau quanh phòng.\\
  Thỏ xanh lùi bước rụng rời,\\
  Chàng ôm mặt, thở chẳng lời nào hơn.\\
  Sấm Âm ngẩng mặt thơ ngây,\\
  Cười tươi rói giữa vòng tay bạn bè:

  \begin{figure}[H]
    \centering
    \begin{subfigure}[t]{0.45\textwidth}
      \centering
      \includegraphics[width=8cm]{../../../../Image/Book?/10-6g1.png}
    \end{subfigure}
    \quad
    \begin{subfigure}[t]{0.45\textwidth}
      \centering
      \includegraphics[width=8cm]{../../../../Image/Book?/10-6g2.png}
    \end{subfigure}
  \end{figure}


  ``Chào chiều, nắng xuống ngoài kia,''\\
  Một giấc trưa đẹp, chẳng chê điểm nào!''

\end{center}

%==========================================TRUYỆN PHỤ=========================================
\omk

\begin{center}
  Thấy mọi người mặt tái xao,\\
  Nàng ngơ ngác hỏi: ``Sao nào, mọi người?''\\
  Thỏ run đáp: ``Chào chị ơi\dots''\\
  Bạn A thì vẫn lặng cười lơ ngơ.\\
  Khuyển xua tay, bảo: ``Dở hơi!\\
  Căng thẳng làm gì, nghỉ ngơi giữa ngày.''\\
  Nói rồi nàng bước nhanh tay,\\
  Tung tăng vào bếp như bay không màng.\\
  Lôi ra chiếc bánh vàng vàng,\\
  Cười vui hí hửng, đặt sang lò nồng.\\
  Chẳng ngờ đĩa lại bằng đồng,\\
  Viễn và A đã kêu không kịp lời.\\
  Chiếc lò nổ phát một hơi,\\
  Như điều hiển hiện giữa đời phàm gian!

  Xung quanh giờ khói lan tràn,\\
  Gạch tường văng vãi từng hàng ngang rơi.\\
  Phù Thủy, Thiên Sứ tơi bời,\\
  Nằm im không thốt một lời nào ra.\\
  Riêng chàng Cửu Viễn tỉnh mà,\\
  Thở dài bất lực nhìn qua Khuyển nằm.\\
  ``Ngủ trưa thì cũng tốt lắm,\\
  Nhưng ai đền khoản chi lắm thế này?''\\
  Khuyển nhoẻn miệng cười lắt lay,\\
  Lẩm bẩm: ``Em ngủ mơ ngày mà thôi\dots''\\

  \img{10-6h}

  May sao chẳng phải mất toi,\\
  Nhờ vào bảo bối cứu đời Viễn mang.\\
  Cả căn phòng trở lại sang,\\
  Mọi tường, mọi góc đều vàng như xưa.\\
  Khuyển vui, cúi tạ bẩm thưa:\\
  ``Cảm ơn anh đã cứu bừa lần này!''\\
  Tuy nhiên vẫn bị mắng thay,\\
  Mụ tóc tím kính, chẳng ai bênh nàng.\\
  Viễn cười mếu, mắt ngỡ ngàng:\\
  ``Chừng nào em mới thôi mang họa về!?''

  \il

  Bạn A đẩy kính quay về,\\
  Thở dài một tiếng: ``Dạo này căng ghê\dots''\\
  Khuyển gãi đầu, cười le te:\\
  ``Tại em hay ngủ mê mê đấy mà\dots''\\
  Viễn chen lời: ``Ừ đúng ha\dots\\
  Còn hay bày chuyện ba hoa chẳng ngừng.''\\
  Mọi người cười phá tưng bừng,\\
  Căn phòng lại sáng tưng bừng như xưa\dots\\
\end{center}

\end{document}